\begin{solucion}
(b) Sea $L$ el cuerpo de descomposición de $f$ sobre $\mathbb{Q}$, es decir, el menor campo que contiene todas las raíces complejas de $f(x)$ [Clase 12, pág. 4]. Por lo que para calcular $[L:\mathbb{Q}]$ observemos que:
\begin{enumerate}
    \item \textbf{$f(x)$ es irreducible en $\mathbb Q[x]$}:Dado que el $1<\partial f=3\leq 3$ entonces $f$ es reducible si y solo si $f$ tiene una raíz. Dado que $f(x)$ no posee raíces racionales, pues por el criterio de la raíz racional las únicas posibles raíces son $\pm1$ y:
    \begin{align*}
        f(1) &= 1^3 - 3\cdot1 + 1 = -1 \\
        f(-1) &= (-1)^3 - 3\cdot(-1) + 1 = 3~.
    \end{align*}
    Por lo tanto, $f(x)$ no tiene raíces racionales, por lo tanto $f$ es irreducible en $\mathbb{Q}[x]$ 
    % (por ejemplo, $f(0)=1$, $f(1)=-1$, $f(-1)=3$
    % [Clase 3, pág. 5], por lo tanto $f(x)$ es irreducible en $\mathbb{Q}[x]$. Dado que $f$ es de grado 3, la única forma en que podría factorizarse sobre $\mathbb{Q}$ sería como el producto de un polinomio lineal y uno cuadrático. La ausencia de raíces racionales descarta la existencia de factores lineales en $\mathbb{Q}[x]$, luego $f$ no es reducible en $\mathbb{Q}$ [Clase 4, pág. 2]. Concluimos que $f(x)$ es un polinomio irreducible sobre $\mathbb{Q}$.
	\item \textbf{f(x) es minimal de $\alpha$ sobre $\mathbb Q$}:  Dado que $f$ es irreducible y $\alpha$ es una de sus raíces, $f(x)$ coincide con el polinomio mínimo de $\alpha$ sobre $\mathbb{Q}$, pues en otro caso si $m(x)$ fuera el minimal de $\alpha$ con $\partial(m)\leq\partial(f)$ entonces $m|f$ no cual contradice que $f$ es irreducible, asi $f$ y $m$ pueden diferir por una constante, pero como $f$ es monico entonces $f=m$. 
    
    \item \textbf{El grado de la extension $[\mathbb Q(\alpha):\mathbb Q]$ es 3} Dado que $f(x)$ es un polinomio minimal de $\alpha$ sobre $\mathbb{Q}$ de grado $3$, y como el grado de una extensión simple generada por un elemento algebraico es igual al grado del polinomio mínimo de dicho elemento entonces se tiene que $[\mathbb{Q}(\alpha):\mathbb{Q}] = \deg(f) = 3$.

	\item \textbf{Campo que contiene a todas las raíces}:  Por el literal (a), sabemos que si $\alpha$ es una raíz de $f$, entonces $\beta = \alpha^2 - 2$ y $\gamma = (\alpha^2-2)^2 - 2$ son las otras dos raíces del polinomio (todas distintas entre sí). Observemos que tanto $\beta$ como $\gamma$ pertenecen a $\mathbb{Q}(\alpha)$, ya que
	\begin{align*}
        \beta = \alpha^2 - 2 \in \mathbb{Q}(\alpha), \qquad \gamma = (\alpha^2 - 2)^2 - 2 \in \mathbb{Q}(\alpha),
    \end{align*}
    Es decir que $\beta$ y $\gamma$ son expresiones racionales de $\alpha$ y por lo tanto pertenecen a $\mathbb{Q}(\alpha)$.
\end{enumerate}

De $4$ se concluye que el cuerpo de descomposición $L$ de $f(x)$ es $\mathbb{Q}(\alpha)$, pues en este el polinomio $f(x)$ se descompone totalmente. Y de $3$ se concluye que el grado de la extensión de descomposición de $f(x)$ sobre $\mathbb{Q}$ es $3$. Por lo tanto, el grado del cuerpo de descomposición $[L:\mathbb{Q}] = [\mathbb{Q}(\alpha):\mathbb{Q}] = 3$.

% por ser combinaciones polinómicas de $\alpha$ con coeficientes en $\mathbb{Q}$. En consecuencia, todas las raíces ${\alpha,\beta,\gamma}$ están contenidas en el campo $\mathbb{Q}(\alpha)$. Esto demuestra que $\mathbb{Q}(\alpha)$ ya es un campo que contiene a todas las raíces de $f(x)$, y por la minimalidad del cuerpo de descomposición se sigue que
% \begin{align*}
%     L = \mathbb{Q}(\alpha)~.
% \end{align*}

% Finalmente, puesto que $L = \mathbb{Q}(\alpha)$, se tiene inmediatamente que
% \begin{align*}
%     [L : \mathbb{Q}] \;=\; [\mathbb{Q}(\alpha) : \mathbb{Q}] \;=\; 3~,
% \end{align*}
% como determinamos en el paso 2. Por lo tanto, el grado de la extensión de descomposición de $f(x)$ sobre $\mathbb{Q}$ es $\mathbf{3}$.


% Conceptos y resultados utilizados:

% \begin{tabular}{ll}
% • Polinomio irreducible: & Definición de polinomio irreducible en $\mathbb{Q}[x]$ (Clase 3, pág. 5).\\
% • Criterio de la raíz racional: & Un polinomio con coeficientes enteros sin raíz racional de la forma $\frac{p}{q}$ 
% es irreducible en $\mathbb{Q}[x]$ si $\deg=3$ (Clase 4, pág. 2).\\
% • Polinomio mínimo: & Existencia/definición del polinomio mínimo de un elemento algebraico (Clase 5, pág. 3).\\
% • Grado de extensión: & $[\mathbb{Q}(\alpha):\mathbb{Q}] = \deg(\text{polinomio mínimo de }\alpha)$ (Clase 6, pág. 2).\\
% • Raíces conjugadas: & Raíces algebraicas que comparten el mismo polinomio mínimo irreducible (Clase 13, pág. 2).\\
% • Cuerpo de descomposición: & Campo más pequeño que contiene todas las raíces de $f(x)$ (Clase 12, pág. 4).\\
% \end{tabular}
\end{solucion}